% !Mode:: "TeX:UTF-8"
% !TEX program  = xelatex

\section{总结与展望}
\label{sec:conclusion}

\subsection{全文总结}

本文研究 TSP 与 CVRP 的算子选择问题。
在比较多种决策粒度设定后，
本文最终将该任务建模为只选初始化的上下文老虎机问题，并实现了 Neural-LinUCB 选择模型。
在方法上，核心设计包括共享动作空间与可行动作掩码、
图编码与手工特征及方法嵌入、按问题拆分的浅层线性 UCB 头，
以及基于历史老虎机反馈的周期性表示层训练。

实验结果表明：
\begin{itemize}
  \item 在合成混合数据集上，Neural-LinUCB 取得最低 mean length，并优于 SingleBest；
  \item 在 LIB 基准数据集上，Neural-LinUCB 的泛化能力有待提升。
\end{itemize}

\subsection{研究局限}
\begin{itemize}
  \item 当前 reward 采用实例内排名映射，其与最终 mean length 目标并不完全一致。
  \item 当前工作只聚焦初始化阶段，没有扩展到“初始化+选择”的两阶段组合。
  \item 当前工作仅覆盖 TSP 和 CVRP 两类问题。
\end{itemize}

\subsection{未来工作}

后续工作可以从以下几个方向继续推进。
\begin{itemize}
  \item 重新设计更贴近最终成本目标的 reward
  \item 扩展到更多组合优化问题类型，验证方法的通用性
  \item 在保持初始化任务清晰可解释的前提下，探索初始化选择与迭代选择的两阶段控制框架。
  \item 开展系统的消融实验，量化图编码器、手工特征、方法嵌入和按问题拆分的UCB 头等各组件对选择性能的边际贡献，以明确模型中各设计选择的必要性。
  \item 扩充基线方法集合，加入更多上下文老虎机变体（如 $\varepsilon$-greedy、Thompson Sampling）以及监督学习基线，以更全面地评估不同决策框架的优劣。
  \item 对关键超参数（如探索系数 $\alpha$、表示层更新周期、嵌入维度等）进行系统的敏感性分析，量化其对选择性能的影响。
\end{itemize}

当前的结果说明了，方法选择可以作为一个独立的研究问题成立。
Neural-LinUCB 已经验证了这一路径的可行性，也为后续更细致的两阶段选择框架打下了基础。
